\chapter*{General Conclusion}\label{chap:gen_con}
\addcontentsline{toc}{chapter}{General Conclusion}

This thesis asked whether quantum and quantum-inspired molecular representations improve antimalarial candidate prioritization over strong classical baselines, and what a genuine quantum machine learning implementation would have to demonstrate on the same chemical space before any claim could be made. The question was addressed through five studies that share one library, one label source and one evaluation protocol: a chemical-space construction and target-anchored docking study, a resistance-aware scoring study stress-tested by explicit-solvent molecular dynamics and bounded by retrospective experimental data, a benchmark of topological, tensor-compressed and quantum-inspired descriptors, a benchmark of graph neural networks and a molecular transformer under chemical distribution shift, and a quantum machine learning programme, refocused on structure-direct encoding, whose protocol is specified and frozen while its canonical benchmark remains pending. The organising result is negative and was stated in advance so that it could not be mistaken for a late qualification: across every representation tested, an extended-connectivity fingerprint with a random forest remained the strongest descriptor for activity prediction on the panels studied, and no fusion, kernel or learned architecture improved on it.

The chemical-space study established the substrate on which every other result rests. An expansion of \num{396} African natural products and \num{454} synthetic antimalarials through complementary similarity-search and SELFIES-mutation routes produced a hybrid library of \num{65856} unique molecules that is \qty{92.6}{\percent} whole-molecule-novel relative to its seeds while retaining \qty{69.3}{\percent} of their annotated scaffolds, an asymmetry that locates the novelty in peripheral substitution rather than in ring systems. Centroid-based sampling reduced the docking campaign from \num{263424} to \num{1936} calculations, and target-anchored docking against four mechanistically distinct proteins, with the chloroquine-resistance transporter corrected to a wild-type receptor and audited by an independent pocket-detection method, resolved the \num{17}-member cohort into seven A*, four B and six C resistance-score classes on the corrected channel. Two properties of that result matter beyond the cohort. First, target breadth and mutation resilience proved separable, because five candidates were classified from a single target panel while remaining non-binders elsewhere. Second, the corrected transporter channel behaved as a null under its own mutation pipeline (\numrange{99.4}{100.5}\%), so the class boundaries were anchored to that null rather than to raw thresholds alone. The docking protocol was validated in parts and bounded honestly: redocking reproduced four of five crystal poses, enrichment on known actives was strong under the dual-filter consensus, and the external benchmark returned a near-chance area under the ROC curve that was retained in the record rather than repaired, with a controlled two-arm experiment excluding cofactor blockade as its cause.

The resistance-aware study converted the docking campaign into a triage framework and then subjected it to the stress tests that determine what it may claim. Static docking and short explicit-solvent dynamics disagreed in a consistent direction: across the eight matched mutant states with an eligible wild-type baseline, docking predicted reduced score retention while the trajectories retained the pose in \num{7} of \num{8} comparisons (\qty{87.5}{\percent}; 95\% Wilson interval \numrange{52.9}{97.8}\%). The disagreement is not an error of either method but a divergence of estimands, and it is used as a calibration flag that prioritizes states for experimental follow-up. Two retrospective checks bounded the static estimand from the experimental side. A curated genotype-resolved inhibition panel showed that rigid-receptor score differences fail to track published fold-shifts at every dihydrofolate-reductase genotype, with direction agreement of only \numrange{28.6}{41.7}\% and the clinically strongest resistance signature mispredicted in direction; and a control on five approved antimalarials failed to recover the documented resistance phenotypes of pyrimethamine and chloroquine, reproducing the null character of the corrected channels on clinically familiar ligands. Together these checks support one operational conclusion: resistance-resilience classes are score-retention categories, useful as triage inputs and as generators of testable hypotheses, and never as resistance measurements.

The two representation benchmarks produced the thesis's central negative result with its mechanisms attached. The descriptor benchmark showed that the persistent-homology fingerprint and the tensor-compressed embedding carry genuine structural information, the former resolving the novelty-versus-scaffold paradox into conserved ring topology and diverging connectivity, the latter retaining docking-relevant information on its derivation panel, yet neither improved activity prediction: the stacked ensemble that included ECFP4 changed its area under the ROC curve by \num{-0.0002}, and the quantum kernel, although it matched a bandwidth-tuned classical kernel (\num{0.823} against \num{0.829}, $p = \num{0.060}$), achieved a fourfold higher kernel--target alignment without converting it into performance, and fell below the same comparator under a label-source shift. The learned-representation benchmark extended the comparison to architectures that learn their own representations. Under random splitting every model sat within a few points of the fingerprint baseline; under scaffold separation the ranking was preserved but the deficits became statistically resolvable, with the graph isomorphism network at \num{0.8047} and the transformer at \num{0.7867} against \num{0.8300} for ECFP4--RF, and a capacity sweep confirmed that the gap is architectural rather than a tuning artifact. The deficit localized in the structural-extrapolation tail and narrowed, without reversing, in the interpolation regime, which converted model choice into a spatial decision and motivated a distance-aware triage rule that remains, by declared design, a testable operational hypothesis rather than a validated procedure. Topological fusion improved the graph network in every structural cohort, but the gain did not scale with structural complexity, its polycyclic interval included zero, and the descriptor salience that accompanied it was not evidence of contribution. Because the activity labels are model predictions on a prevalence-enriched panel, all early-recognition metrics saturate at their prevalence-imposed ceilings, and no prospective screening precision is claimed from any of these results.

These findings are bounded by limitations that are stated rather than smoothed. Docking scores are scoring-function outputs on anchored poses, not affinities; the transporter and ion-pump sites are proxy or mechanism-motivated cavities; the molecular-dynamics evidence is single-replicate at \qty{10}{\ns} and supports pose-retention statements only, with endpoint energies carrying an inter-replicate sensitivity of up to \qty{7.75}{\kcalmol}; the activity labels are predictions from an external model on a scaffold-poor panel, so every classification metric measures agreement with one model's opinion of activity; the resistance cohorts are small, and at $n = \num{17}$ a correlation test reaches \qty{80}{\percent} power only above a moderate association, so every null association is reported as a failure to detect rather than as evidence of equivalence; and no experimental validation, biochemical or phenotypic, was performed within this work, which is why every candidate shortlist is presented as a hypothesis for assay follow-up rather than as a hit list. The honest-negative framing is itself part of the contribution: each null result is documented with the analysis that would have detected the opposite, so that the boundary is auditable rather than asserted.

The forward direction is the quantum machine learning programme, whose central question the completed studies made precise. The benchmarks exposed two structural blind spots in the representations examined: hashed fingerprints collapse stereochemistry, and local message passing loses ring geometry across unseen scaffolds; and the descriptor benchmark showed that a quantum component fed by an ECFP4-derived input inherits the fingerprint's information ceiling, which is why its favourable Hilbert-space geometry produced no performance difference. The programme was therefore refocused on quantum molecular structure encoding, which places bond-order and Coulomb-adjacency information directly into one- and two-qubit rotation angles, carries stereocentre configuration and double-bond geometry as signed parameters, and trains a hybrid quantum-classical model end to end \cite{Boy2025QMSE}. Its encoding substrate is frozen on the complete \num{19849}-molecule panel with zero failures, its training pipeline is confirmed end to end, and its canonical benchmark, defined on the same frozen partitions, label provenance and multiplicity machinery as every result above, is pending; three outcome scenarios are predeclared, spanning advantage, equivalence and a depth-limited deficit, and each is interpretable under the claim-ledger convention of this thesis. Hardware evaluation follows only after the simulator benchmark is complete under that protocol.

Taken together, the thesis delivers three things to the field. It provides a fully declared evaluation substrate, with frozen panels, partition families, power boundaries and a claim ledger, that any future quantum implementation on this chemical space inherits rather than rebuilds. It provides a reusable resistance-aware workflow whose components, the wild-type eligibility gate, the estimand-divergence diagnostic and the retrospective phenotype confrontation, transfer to other resistance-relevant target panels, a transfer whose timeliness grows with the documented expansion of artemisinin-resistance markers across Africa. And it provides a documented record of what does not work: novel representations, whether topological, tensor-compressed, kernel-based, graph-based or sequence-based, add diagnostic value without predictive gain on these panels, and the credibility of that record rests on the machinery that would have detected the opposite. The next step is the one the thesis was built to enable: a structure-direct quantum benchmark on this library, judged against the comparison standard the completed studies define.
